% mage-auto-heading: engineering-through-representations
% mage-auto-heading: why-software-often-modeled-differently
% mage-auto-heading: model-based-engineering-at-smaller-scales
% mage-auto-heading: model-first-engineering-as-one-point-in-the-methodological-spectrum
% Authoritative LaTeX manuscript; Markdown is a generated compatibility projection.

% <!-- part-title: The Profession -->

% <!-- chapter-title: Back to Models, but Not Back to Waterfall -->

\chapter{Back to Models, but Not Back to Waterfall}

MAGE asks engineers to invest deliberately in models, specifications, and controls while delegating more implementation to agents. Does that mean software engineering is returning to waterfall: specifying the system in advance, then handing the specification to a fabricator? The resemblance is understandable. Much of software engineering\textquotesingle s history has involved moving away from the expectation that a sufficiently complete design could precede implementation. Yet the iterative development examined in \magevalue{chapter:one-problem-many-models} showed the opposite of a fixed handoff: implementation exposed missing knowledge, models were revised, and recurring failures changed the engineering environment inherited by later work.

The distinction is between working through models and requiring models to settle the implementation before work begins. Engineers have long used representations to make consequential decisions without expecting those representations to describe every detail of the finished artifact. Software engineering developed a different balance because code itself served as an unusually powerful representation and because maintaining additional models imposed costs. Commodity intelligence changes those costs. To understand what MAGE changes, and what it inherits, we need to examine why software engineering historically organized representation and implementation differently from other engineering disciplines.

\section{Engineering Through Representations}\label{heading:engineering-through-representations}

% <!-- index-def: universal-language -->

The account of Modeling in \magevalue{chapter:modeling} belongs to a much older engineering practice. Civil, mechanical, electrical, aerospace, and systems engineers work through load models, schematics, state-space models, simulations, drawings, and analyses. Models are the building blocks of engineering.\mageaside{Here I mean model in the broad sense developed in \magevalue{chapter:modeling}: a purposeful reduction that preserves distinctions needed for some question. The claim is close to axiomatic under this book\textquotesingle s definition of Engineering. A consequential system sufficiently simple to understand and direct without reasoning through reduced representations does not present the complexity for which engineering, in this sense, is required.} A circuit schematic is not a circuit; a finite-element model is not a bridge; a process-flow diagram is not a refinery. Each representation preserves selected properties of the system while omitting others. In physical engineering, where fabrication can be expensive or irreversible, these representations allow engineers to evaluate designs before committing to realization. A structural model can answer a loading question before concrete is poured; circuit simulation can expose a design error before a mask set is fabricated.

\textbf{The common engineering move is simple: ask the question over the cheapest representation that preserves the property you care about, then use evidence to justify realization.} Mature engineering still uses prototypes, experiments, and direct inspection. Models do not replace the artifact. They organize the reasoning around it.

% <!-- point: software-makes-the-same-move-through-its-own-traditions | Requirements engineering, architecture, type systems and intermediate representations, formal methods, and model-driven approaches are software's own manifestations of one engineering move: represent the property at a level where it can be reasoned about. -->

Software makes the same move through its own traditions. Requirements engineering externalizes obligations. Architecture exposes consequential structural relationships. Type systems and intermediate representations make selected program properties explicit. Formal methods construct representations over which stronger claims can be established. Model-driven and model-based approaches go further and make selected models primary engineering artifacts. These traditions differ substantially, and they share one move: represent the property at a level where it can be reasoned about.

The medium varies; the purpose does not. The analysis serves a decision. Salado makes this point directly by treating engineering as disciplined decision-making under uncertainty and competing objectives, rather than technical problem solving alone~\citep{salado2026}. Engineers model, calculate, measure, and experiment because consequential choices must be made: which design to pursue, what evidence to gather, which tradeoff to accept, or whether the resulting system is fit to proceed. A model earns its place when it stands in service of such a decision.

\section{Why Software Often Modeled Differently}\label{heading:why-software-often-modeled-differently}

Software engineering did not reject models in favor of iteration. Its economics made implementation itself an unusually attractive surface for engineering reasoning.

Software engineering already combines explicit obligations with rapid iteration. Test-driven development, for example, states an expected behavior in a test before implementing it, then revises the implementation through repeated test–code–refactor cycles~\citep{beck2003tdd}. Specifying an obligation before implementing it is not the same as specifying the entire system before development begins. Likewise, agile methods do not dispense with planning. They favor delivering useful results early, learning from them, and revising subsequent plans as evidence accumulates~\citep{agile-principles2001}.

\textbf{Software realization was extraordinarily cheap; detailed design and implementation labor were not.} Once source code existed, another compilation, execution, or copy of the program could usually be produced at negligible marginal cost. Producing, understanding, debugging, and changing that source consumed skilled engineering labor. The source was already a precise, executable representation through which engineers reasoned about the system.

% <!-- point: code-carried-representation-and-realization-at-once | Software design and implementation have historically been tightly coupled: the artifact handed to the compiler is itself the product of an iterative design process, so code served as realization, executable specification, and the primary surface for further reasoning. -->

That blurred a distinction the older disciplines could make more cleanly. A structural engineer reasons through a design before concrete is poured; a semiconductor designer simulates a circuit before fabrication. In software, the artifact handed to the cheap realizing mechanism, the compiler, is itself the product of an iterative design process. Reeves made the strongest version of this argument by calling source code the final software design. MAGE does not require that full claim. It requires only the narrower one: \textbf{software design and implementation have historically been tightly coupled, and code has served simultaneously as realization, executable specification, and a primary surface for further reasoning.}\mageaside{Jack W. Reeves\textquotesingle s "code is design" argument~\citep{reeves1992whatissoftwaredesign} is a useful historical extreme; MAGE relies only on the narrower claim that software design and implementation have traditionally been tightly coupled.}

Code\textquotesingle s malleability then works against the engineer as a system evolves. Architecture, intent, constraints, and operational knowledge can stay implicit in an artifact whose structure keeps changing, and the system stays executable long after the knowledge needed to govern it has become hard to recover. The disciplines named in the previous section exist partly to hold that knowledge where continual change cannot quietly erase it.

Those secondary models faced an unusual economic test. The implementation was already precise, executable, and immediately available, while an architectural, behavioral, or process model imposed an additional correspondence burden. A model could be useful when authored and misleading after the implementation moved beyond the relation it claimed. Many projects therefore kept such representations thin unless their communication, reasoning, or assurance value repaid that standing synchronization cost.

% <!-- point: commodity-intelligence-moves-the-price-ratio | Commodity intelligence lowers the labor of realizing software and parts of the labor of deriving, updating, reconciling, and checking additional representations, so implementation and live modeling become cheaper together and model-mediated software engineering reaches surfaces the synchronization tax once excluded. -->

Commodity intelligence changes both sides of that ledger. It lowers the skilled labor required to realize and revise software, and parts of the labor required to derive, update, reconcile, and check additional representations. \textbf{Implementation and live modeling can become cheaper together}, making model-mediated software engineering economical over surfaces where the synchronization tax once exceeded its return. The construction and carrying costs examined in \magevalue{sec:scope-conditions} still determine where that investment pays.

% <!-- point: the-separation-is-not-absolute | Deployed software determines realized behavior; models can express intent, constrain implementation, or supply executable abstractions without displacing implementation as an engineering artifact. -->

Implementation remains the final authority on what the software actually does. Models may describe intended behavior, constrain implementation, or supply executable abstractions, but the deployed system determines its realized behavior. Some properties are therefore most naturally reasoned about through implementation itself. What changes is that code need no longer bear the whole burden of being the realization, the repository of system knowledge, and the surface over which engineering reasoning occurs.

\section{Model-Based Engineering at Smaller Scales}\label{heading:model-based-engineering-at-smaller-scales}

Model-based engineering has long involved investment in the machinery that makes models useful. An established environment can provide modeling languages and editors, analyses, simulation, transformations, and verification tools. Engineers can then extend models within the supported vocabulary without constructing new machinery for every additional element. The investment can serve many engineering activities and, where the domain permits reuse, several products~\citep{schmidt2006model,kelly2008domain}.

Domain-specific modeling makes this economic relationship explicit. A vocabulary suited to a domain and its associated generators and analyses can reduce repeated work, but constructing them requires knowledge of both the domain and the tooling. Maintenance adds to that investment as the language, platform, and engineering needs change~\citep{kelly2008domain,kautz2018mbse}. Whether the arrangement pays therefore depends on the problems it serves and how often the machinery remains useful.

Commodity intelligence may change the scale at which such investments become economical. If agents reduce the effort required to construct and maintain a representation and its consumers, a project, subsystem, or individual engineering problem may justify machinery that previously required broader reuse to repay its cost. An organization with established model-based engineering practices may instead use agents to extend its existing assets. In either case, engineers must still decide which questions warrant models and whether the resulting machinery supplies dependable evidence. Lower construction costs expand the available choices; they do not settle those decisions.

\section{Model-First Engineering as One Point in the Methodological Spectrum}\label{heading:model-first-engineering-as-one-point-in-the-methodological-spectrum}

If engineers can economically introduce modeling machinery for individual engineering problems, when should they construct the models themselves? At one end of the spectrum, \emph{model-first} engineering develops representations of the intended system before realization, using them to guide design, analysis, and implementation. At the other, \emph{model-last} engineering begins with a working implementation and develops models afterward, to understand its behavior, recover its architecture, or establish properties that have become consequential. These are two possible arrangements of the engineering work, not mutually exclusive commitments.

The Siemens account in \autoref{chapter:5.4-other-agentic-software-factories} illustrates a model-first arrangement in an established systems-engineering environment. Agents work through SysML v2 models and engineering simulations to explore candidate designs and evaluate them against requirements~\citep{liu2024agentstudio}. The models are already part of the organization\textquotesingle s engineering practice; agents can use them rather than reconstructing the system from realized artifacts. Siemens thus provides a concrete example of agents entering a production process organized around engineering representations.

% Evidence: E01 | source: book/data/local-evidence/evidentiary-basis-11202026/integration/records/E01.md | kind: author-account
% Evidence: E25 | source: book/data/local-evidence/evidentiary-basis-11202026/integration/records/E25.md | kind: author-account
% Evidence: E26 | source: book/data/local-evidence/evidentiary-basis-11202026/integration/records/E26.md | kind: dated-case-evidence
Consider the advantages Siemens enjoys from decades of investment in model-based engineering~\citep{mathieu2020amesim}. Its agents enter an environment already equipped with engineering representations, analyses, simulations, and tools for working with them. But consider also what I have described in this book. I now have a model-based engineering environment of my own, suited to the particular problems of DocAble. I built it alongside the product, without a dedicated tooling team, in a matter of months rather than decades. As a result, I can enjoy many of the same benefits. I can explore alternative designs through models, investigate the consequences of proposed changes before implementing them, and delegate substantial implementation work with confidence that the obligations I have modeled and made enforceable will continue to govern that work. This may be one of the more important consequences of commodity intelligence for engineering practice: engineers need not confine themselves to the modeling languages and tools already available to them. They can increasingly afford to construct their own.

Several other factories in \magevalue{chapter:software-factories} start from existing software. Cloudflare, Shopify, and GitLab introduce agents into established software-development environments organized around source repositories, tests, and operational controls~\citep{reimann2026codex, moreno2026river, gitlab2026customize}. The software already exists, and agents must work through or recover the engineering knowledge embedded in it. These implementation-centered starting points can support model-last work, although the accounts do not establish that models followed implementation or describe a uniform methodology for constructing and maintaining them.

% Evidence: E25 | source: book/data/local-evidence/evidentiary-basis-11202026/integration/records/E25.md | kind: author-account
% Evidence: E26 | source: book/data/local-evidence/evidentiary-basis-11202026/integration/records/E26.md | kind: dated-case-evidence
A project need not occupy a fixed point on this spectrum. \magevalue{chapter:one-problem-many-models} traced how DocAble\textquotesingle s engineering representations developed alongside implementation. Some models were constructed before particular changes to guide realization; others emerged when implementation exposed missing abstractions, recurring failures, or obligations worth enforcing. Test-driven development provides a familiar precedent: expected behavior is specified before implementation, but tests and implementation evolve together through repeated cycles~\citep{beck2003tdd}. The important requirement is not that every model precede implementation, but that consequential properties, design decisions, and their rationales become explicit before an implementation is accepted as satisfactory.

% <!-- point: mage-includes-model-first-organization-without-privileging-it | Degrees of Freedom and Tolerance help engineers decide which choices to settle through models and which to leave to realization; the balance depends on the engineering question and the costs and benefits of available representations. | terms: degrees-of-freedom, tolerance -->

MAGE uses the concepts of Degrees of Freedom and Tolerance to help engineers decide which choices to settle through models and which to leave to realization. Where models determine the realization sufficiently completely, deterministic generation may be preferable. Where consequential choices remain open, engineers can preserve the obligations and tolerances that matter while leaving agents to choose among acceptable implementations. Those degrees of freedom are not defects in the model; they are choices engineers have not, or could not, settle in advance. The appropriate balance depends on the engineering question and the costs and benefits of the available representations, potentially differing even among components or properties of the same system.
